\chapter{Gestão do Projeto}

A gestão do projeto foi conduzida de forma colaborativa, considerando a divisão de responsabilidades entre os integrantes da equipe, o acompanhamento contínuo das atividades e a necessidade de cumprir os prazos definidos para as entregas acadêmicas. Por se tratar de um projeto desenvolvido em equipe e com diferentes frentes técnicas, como front-end, back-end, banco de dados e documentação, a organização das tarefas foi essencial para manter a evolução do sistema e garantir a integração entre as partes.

Foram adotados princípios de metodologias ágeis, adaptados ao contexto acadêmico e à disponibilidade dos integrantes. As atividades foram distribuídas em ciclos de desenvolvimento, com definição de prioridades, acompanhamento das entregas e revisão contínua das funcionalidades implementadas. Embora o processo não tenha seguido rigorosamente todos os rituais formais de um framework ágil, a equipe utilizou práticas como divisão de tarefas, acompanhamento por quadro Kanban, reuniões de alinhamento e versionamento contínuo do código.

A gestão também envolveu o uso de ferramentas colaborativas para comunicação, documentação, controle de versão e acompanhamento das atividades. Esses recursos contribuíram para organizar o fluxo de trabalho, registrar decisões importantes e facilitar a integração entre os membros da equipe durante o desenvolvimento do sistema Tropical Mix.

\section{Organização da Equipe}

A equipe do projeto foi composta por cinco integrantes, com responsabilidades distribuídas de acordo com as áreas técnicas envolvidas no desenvolvimento do sistema. A divisão considerou as principais frentes do projeto, incluindo gerenciamento das atividades, desenvolvimento do front-end, implementação do back-end, modelagem do banco de dados, documentação e integração entre as camadas da aplicação.

Apesar da divisão por áreas, o desenvolvimento ocorreu de forma colaborativa, com revisões conjuntas, apoio entre os membros e alinhamentos frequentes para garantir a compatibilidade entre as partes implementadas. Essa organização foi importante para manter a continuidade das entregas e facilitar a identificação de dependências entre interface, API e banco de dados.

A Tabela~\ref{tab:integrantes-projeto} apresenta os integrantes da equipe, suas funções principais e as responsabilidades associadas a cada atuação no projeto.

\begin{table}[H]
\centering
\caption{Integrantes do projeto e responsabilidades.}
\label{tab:integrantes-projeto}
\resizebox{\textwidth}{!}{%
\begin{tabular}{|p{4cm}|p{4cm}|p{8cm}|}
\hline
\textbf{Integrante} & \textbf{Função principal} & \textbf{Responsabilidades e tecnologias utilizadas} \\ \hline

Gabriella Santos Mendes & 
Gerente de Projeto / Desenvolvedora Front-end & 
Planejamento e acompanhamento das atividades; desenvolvimento da interface do usuário utilizando JavaScript, React, HTML e CSS; integração com o back-end. \\ \hline

Gustavo Almeida de Oliveira & 
Desenvolvedor Back-end & 
Apoio no desenvolvimento da API em Python com Flask; estruturação das funcionalidades e testes do servidor. \\ \hline

Isabelle Silva Severino & 
Administradora de Banco de Dados (DBA) & 
Modelagem e implementação do banco de dados MySQL; criação de tabelas, relacionamentos e consultas. \\ \hline

Matheus Fogaça Sarretto & 
Desenvolvedor Back-end & 
Implementação da lógica de negócio e das rotas da API utilizando Python com Flask; integração com o banco de dados. \\ \hline

Matheus Victor Feliciano Jupi & 
Desenvolvedor Front-end & 
Desenvolvimento de componentes e páginas em React, com uso de HTML e CSS; implementação de funcionalidades visuais e integração com a API. \\ \hline

\end{tabular}%
}
\fonte{Elaborado pelos autores.}
\end{table}


\section{Metodologias de Gestão e Desenvolvimento}

O desenvolvimento do sistema Tropical Mix foi conduzido com base em princípios de metodologias ágeis, adaptados ao contexto acadêmico, à disponibilidade dos integrantes e aos prazos estabelecidos para o projeto. A equipe utilizou práticas como divisão de tarefas, definição de prioridades, acompanhamento contínuo das entregas e reuniões de alinhamento para orientar o avanço das funcionalidades.

Embora o processo não tenha seguido integralmente todos os rituais formais de uma metodologia ágil específica, foram adotados elementos inspirados no Scrum e no Kanban. O planejamento foi organizado em ciclos de desenvolvimento, denominados sprints, com objetivos definidos para cada período. Já o acompanhamento das tarefas foi realizado com apoio do GitHub Projects, por meio de um quadro Kanban que permitiu visualizar atividades pendentes, em andamento e concluídas.

Essa abordagem favoreceu a organização do trabalho em equipe, a identificação de dependências entre as áreas técnicas e o acompanhamento do progresso do sistema. Como o projeto envolveu diferentes frentes, como front-end, back-end, banco de dados, documentação e integração, a gestão das atividades foi fundamental para reduzir retrabalho e manter o alinhamento entre os integrantes.

As entregas do projeto foram organizadas em marcos principais, conforme apresentado na Tabela~\ref{tab:marcos-projeto}.

\begin{table}[H]
\centering
\caption{Marcos principais do projeto.}
\label{tab:marcos-projeto}
\resizebox{\textwidth}{!}{%
\begin{tabular}{|p{4cm}|p{3cm}|p{9cm}|}
\hline
\textbf{Marco} & \textbf{Data prevista} & \textbf{Descrição} \\ \hline

Prova de Conceito & 
10/10/2025 & 
Estruturação inicial do sistema, com implementação das tecnologias base e definição da arquitetura. \\ \hline

Entrega da Documentação & 
07/11/2025 & 
Elaboração e padronização dos capítulos do relatório técnico. \\ \hline

Entrega do MVP & 
05/12/2025 & 
Versão funcional do sistema com as principais funcionalidades de controle de estoque e vendas. \\ \hline

Integração e Testes & 
20/04/2026 & 
Validação das funcionalidades integradas do sistema e realização de testes internos e com usuários. \\ \hline

Homologação do Sistema & 
17/05/2026 & 
Correções finais, ajustes de desempenho e homologação da versão estável. \\ \hline

Entrega Final do Projeto & 
27/05/2026 & 
Entrega oficial do projeto completo, incluindo documentação técnica. \\ \hline

\end{tabular}%
}
\fonte{Elaborado pelos autores.}
\end{table}

\subsection{Sprints}

O desenvolvimento do projeto foi organizado em ciclos de trabalho denominados sprints, utilizados como referência para distribuir as atividades ao longo do período de execução. Cada sprint contemplou objetivos específicos, atividades previstas e entregas esperadas, permitindo à equipe acompanhar a evolução do sistema de forma progressiva.

Considerando o contexto acadêmico e a disponibilidade dos integrantes, as sprints foram adaptadas conforme as demandas do projeto, os prazos de entrega e as dependências entre as áreas técnicas. Dessa forma, algumas atividades foram replanejadas ao longo do desenvolvimento, especialmente nos momentos de integração entre front-end, back-end e banco de dados.

O acompanhamento das atividades foi realizado por meio de um quadro Kanban no GitHub Projects, no qual foram registradas tarefas relacionadas à documentação, front-end, back-end, banco de dados, hospedagem, testes e validação. A Tabela~\ref{tab:sprints-projeto} apresenta uma síntese da organização das sprints do projeto, consolidando as principais atividades acompanhadas no quadro.

\begin{table}[H]
\centering
\caption{Sprints do projeto.}
\label{tab:sprints-projeto}
\scriptsize
\resizebox{\textwidth}{!}{%
\begin{tabular}{|p{1.8cm}|p{2.4cm}|p{3.8cm}|p{6.3cm}|p{3.4cm}|}
\hline
\textbf{Sprint} & \textbf{Período} & \textbf{Objetivo geral} & \textbf{Atividades principais} & \textbf{Entregas} \\ \hline

Sprint 1 &
23/09 -- 24/10 &
Planejamento e estruturação inicial do projeto. &
Levantamento do problema; pesquisa de sistemas concorrentes; definição inicial do escopo; criação da documentação do desenho da aplicação; criação do modelo de documentação em LaTeX; configuração inicial dos ambientes de front-end, back-end e banco de dados. &
Prova de conceito, documentação inicial e ambientes configurados. \\ \hline

Sprint 2 &
25/10 -- 21/11 &
Modelagem e construção da base técnica da aplicação. &
Criação do modelo conceitual e lógico do banco de dados; criação do esquema e das tabelas principais; início da documentação do banco de dados; criação dos endpoints iniciais da API; implementação da lógica de registro e login de usuário; criação da estrutura SPA com as páginas necessárias. &
Base do sistema estruturada, banco modelado e primeiras rotas implementadas. \\ \hline

Sprint 3 &
22/11 -- 05/12 &
Implementação do MVP e dos fluxos principais. &
Desenvolvimento das telas de login e cadastro; integração com API de login e registro; criação de componentes do front-end; criação da tela home; desenvolvimento inicial da tela de estoque; implementação das operações de criação, leitura, atualização e exclusão para estoque e vendas no back-end; conexão do front-end com a API. &
MVP funcional com autenticação, telas iniciais, estoque e vendas. \\ \hline

Sprint 4 &
06/12 -- 31/01 &
Integração, hospedagem e expansão das funcionalidades. &
Hospedagem do projeto em ambiente AWS; ajustes da API para a nova hospedagem; alinhamento entre hospedagem e banco de dados; criação da rota para relatório; criação de relatório de estoque ou vendas; povoamento do banco com dados de produtos da loja; ajustes na API e nas integrações. &
Aplicação hospedada, banco populado e integração entre camadas ajustada. \\ \hline

Sprint 5 &
01/02 -- 31/03 &
Refinamento das telas e ampliação dos módulos administrativos. &
Adição de filtros nas telas; implementação de dropdowns; padronização das formas de adicionar produtos; desenvolvimento da tela de fornecedores; ajustes nas tabelas de produto e estoque; ajustes no código e nas rotas; implementação da página de gerenciamento de usuários; ajustes na tabela de usuários. &
Módulos administrativos ampliados, filtros implementados e estrutura de usuários ajustada. \\ \hline

Sprint 6 &
01/04 -- 30/04 &
Aprimoramento do controle de estoque e validação dos fluxos. &
Implementação de regras para não adicionar ou vender produtos fora da validade; criação de view de notificação de validade; desenvolvimento da tela principal de produtos e da tela de detalhes; implementação da tela de clientes; criação do campo de notificações; ajustes de usabilidade e validação dos fluxos principais. &
Controle de validade aprimorado, telas de produtos e clientes implementadas e notificações adicionadas. \\ \hline

Sprint 7 &
01/05 -- 27/05 &
Testes, revisão e documentação final. &
Realização de testes unitários, testes internos e testes com usuário; adição de referências na documentação; revisão da documentação técnica; ajustes finais no front-end, back-end e banco de dados; organização da entrega final do projeto. &
Testes realizados, documentação final em andamento e projeto preparado para entrega. \\ \hline

\end{tabular}%
}
\fonte{Elaborado pelos autores com base no quadro Kanban do projeto.}
\end{table}

\section{Ferramentas de Comunicação e Colaboração}

Para garantir o alinhamento entre os integrantes, foram utilizadas ferramentas colaborativas voltadas à comunicação, controle de versão, documentação e acompanhamento das tarefas. Essas ferramentas auxiliaram na organização do projeto, no registro das decisões e na integração entre as diferentes frentes de desenvolvimento.

A Tabela~\ref{tab:ferramentas-colaboracao} apresenta as principais ferramentas utilizadas pela equipe durante o desenvolvimento do sistema Tropical Mix.

\begin{table}[H]
\centering
\caption{Ferramentas de comunicação e colaboração utilizadas no projeto.}
\label{tab:ferramentas-colaboracao}
\resizebox{\textwidth}{!}{%
\begin{tabular}{|p{4cm}|p{11cm}|}
\hline
\textbf{Ferramenta} & \textbf{Descrição} \\ \hline

GitHub & 
Plataforma utilizada para hospedagem do código-fonte, controle de versão, colaboração entre os integrantes e acompanhamento das tarefas por meio do GitHub Projects. \\ \hline

Google Meet & 
Ferramenta utilizada para reuniões online, alinhamentos quinzenais, compartilhamento de tela e revisão conjunta de funcionalidades e documentos. \\ \hline

WhatsApp & 
Canal utilizado para comunicação rápida entre os integrantes, troca de links, envio de avisos, acompanhamento de pendências e alinhamentos diários. \\ \hline

Reuniões presenciais semanais & 
Encontros realizados durante o horário de aula, com acompanhamento do professor orientador, para validação do andamento do projeto e definição das próximas atividades. \\ \hline

Overleaf & 
Plataforma utilizada para escrita colaborativa da documentação acadêmica em LaTeX, permitindo edição simultânea e padronização do relatório técnico. \\ \hline

Notion & 
Ferramenta utilizada para apoio ao planejamento, registro de anotações, organização de ideias e documentação complementar do projeto. \\ \hline

\end{tabular}%
}
\fonte{Elaborado pelos autores.}
\end{table}

\section{Repositório da Aplicação}

O código-fonte e a documentação do projeto foram armazenados em um repositório público no GitHub, com o objetivo de garantir controle de versão, colaboração entre os integrantes e rastreabilidade das alterações realizadas ao longo do desenvolvimento. O uso do repositório permitiu acompanhar a evolução do sistema, registrar modificações no código e organizar as entregas relacionadas às diferentes frentes do projeto.

O repositório pode ser acessado pelo seguinte endereço:

\begin{center}
\url{https://github.com/GabriellaSMendes/projetoextensao.git}
\end{center}

A estrutura do repositório foi organizada de forma modular, contendo diretórios específicos para as principais partes da aplicação. A pasta \texttt{frontend} concentra os arquivos relacionados à interface web desenvolvida em React. A pasta \texttt{backend} reúne a API desenvolvida em Python com Flask. A pasta \texttt{database} armazena scripts e arquivos relacionados à estrutura do banco de dados. Já a pasta \texttt{LaTeX} contém os arquivos utilizados na elaboração da documentação acadêmica do projeto.

Além dessas pastas principais, o repositório também contém arquivos de configuração, materiais complementares e documentos de apoio ao desenvolvimento. Essa organização contribui para a manutenibilidade do projeto, facilita a localização dos arquivos e permite que os integrantes atuem em diferentes áreas sem comprometer a estrutura geral da aplicação.

O GitHub também foi utilizado como ferramenta de apoio à gestão das tarefas, por meio do recurso GitHub Projects. Esse recurso possibilitou o acompanhamento das atividades em formato de quadro Kanban, permitindo visualizar tarefas pendentes, em andamento e concluídas. Dessa forma, o repositório não atuou apenas como ambiente de armazenamento do código, mas também como instrumento de organização e acompanhamento do desenvolvimento.